\pagenumbering{roman}
\setcounter{page}{1}

\selecthungarian

%----------------------------------------------------------------------------
% Abstract in Hungarian
%----------------------------------------------------------------------------
\chapter*{Kivonat}\addcontentsline{toc}{chapter}{Kivonat}

Jelen dokumentum egy diplomaterv sablon, amely formai keretet ad a BME Villamosmérnöki és Informatikai Karán végző hallgatók által elkészítendő szakdolgozatnak és diplomatervnek. A sablon használata opcionális. Ez a sablon \LaTeX~alapú, a \emph{TeXLive} \TeX-implementációval és a PDF-\LaTeX~fordítóval működőképes.


\vfill
\selectenglish


%----------------------------------------------------------------------------
% Abstract in English
%----------------------------------------------------------------------------
\chapter*{Abstract}\addcontentsline{toc}{chapter}{Abstract}

In the development of safety-critical systems, a significant share of effort and cost lies not in implementation but in validating and verifying the satisfaction of system requirements. In many development processes and industry standards, each requirement must be accompanied by a test strategy that captures its conditions, expected outcomes, and coverage. These artifacts are costly to author and to keep consistent as requirements evolve. As Large Language Models (LLMs) become increasingly used to automate various software engineering tasks, including software testing, they also appear as a promising way to automate safety-critical requirements and test management. However, in safety-critical settings, every generated result must be grounded in and provably consistent with the underlying system model and industry standards, which makes unverified use of LLMs unacceptable due to possible mistakes and logic errors in output, as well as hallucinations. This work presents an approach that couplesLLMs with formal, model-based verification to generate test strategies automatically from natural-language requirements. The generation is grounded in a formal metamodel of the system, and an automated verification loop rejects and refines any result that violates the model's rules.Thus, only consistent strategies are produced. We proposea neuro-symbolic method along with a traceability mechanism that links each generated element back to its originating requirement. We implement our approach in the open-source Refinery graph reasoning framework and demonstrate its applicability and effectiveness with a case study in a safety-critical railway domain inspired by an industrial problem. The framework enables fast yet trustworthy test-strategy creation: engineers only review formally verified output instead of authoring it from scratch, reducing effort while preserving the consistency and correctness of model-based systems. 

\vfill
\cleardoublepage

\selectthesislanguage

\newcounter{romanPage}
\setcounter{romanPage}{\value{page}}
\stepcounter{romanPage}